\section{Roadmap and open risks}
\label{sec:limits}

\subsection{Phases}

\begin{table}[H]
\centering
\begin{tabularx}{\textwidth}{l L{0.46\textwidth} Y}
\toprule
\textbf{Phase} & \textbf{Scope} & \textbf{Status} \\
\midrule
0 & Entry node install; single-machine pool; prove the submit path & \textbf{done} \\
3 & Remote submission: IDTOKENS over one port, no shell anywhere & \textbf{done} (brought forward) \\
--- & Contention-driven owner policy & \textbf{done} \\
1/2 & Second and third workers; daemon tokens; jobs distribute & \textbf{done} \\
--- & Multi-homing across two unroutable subnets & \textbf{done} \\
--- & First user PC: worker \emph{and} submit client & \textbf{done} \\
4 & JupyterHub with \kn{batchspawner} CondorSpawner & not started \\
5 & ARM workers; GPU discovery & not started \\
6 & Overlay VPN for off-LAN machines & deferred by decision \\
\bottomrule
\end{tabularx}
\caption{Phase status.}
\end{table}

Phase 3 was pulled ahead of phases 1 and 2 because it is the requirement the
whole design exists to satisfy. Proving it early de-risked everything after it:
had remote submission turned out to be impractical, the rest of the work would
have been wasted.

\subsection{Phase 4 is blocked on a storage decision}

JupyterHub with \kn{batchspawner} spawns each user's notebook kernel as a Condor
job on a worker, which is exactly the requirement that Jupyter be allowed but
never run on the entry node.

\begin{keypoint}
It is blocked on a decision that is not a Condor question: \textbf{where notebook
home directories live}. Every user needs the same path visible on whichever
worker their kernel lands on. Scratch sandboxes evaporate when the job ends, and
notebooks must persist across sessions and across machines.
\end{keypoint}

That means an NFS export, and it is unprovisioned.

\subsection{Heterogeneity, deferred to phase 5}

\begin{table}[H]
\centering
\begin{tabularx}{\textwidth}{L{0.24\textwidth} Y}
\toprule
\textbf{Item} & \textbf{Detail} \\
\midrule
Architectures & x86\_64 and aarch64 \\
Default safety net & \kn{condor\_submit} pins \kn{Requirements} to the submitting machine's architecture, so ARM capacity is invisible until explicitly opted into \\
Cross-arch jobs & Interpreted code or multi-arch containers only \\
Condor versions & Distribution-supplied versions differ per OS. Pin deliberately; the central manager should not be older than its workers \\
GPUs & \kn{condor\_gpu\_discovery}, \kn{DetectedGPUs}, \kn{request\_gpus}. Jetson iGPU detection under L4T is not reliably clean, so verify rather than assume \\
\bottomrule
\end{tabularx}
\caption{Heterogeneity handling.}
\end{table}

The architecture default is the right one: making ARM capacity opt-in means a
user who has never thought about architecture cannot accidentally match a job to
a machine that cannot run it.

\subsection{Open risks}

\subsubsection{Resolved}

\begin{table}[H]
\centering
\begin{tabularx}{\textwidth}{L{0.36\textwidth} Y}
\toprule
\textbf{Item} & \textbf{Resolution} \\
\midrule
Two unroutable subnets in one pool & Solved via \kn{PRIVATE\_NETWORK\_*} with a derived, not configured, name (Chapter~\ref{sec:mobility}) \\
Role alias for the entry node & Retired. \kn{avahi-publish} is not interface-aware; a machine's own hostname is \\
Policy drift between entry node and workers & Both run the same cron script and the same two-tier policy shape, differing only where the schedd role requires it \\
cgroup per-job enforcement does not work on the installed build & Resolved \emph{by decision}, not by fix. Per-job caps are not wanted on machines this size; \kn{CGROUP\_MEMORY\_LIMIT\_POLICY = none} makes the configuration honest about it \\
Suspend, memory eviction, admission closure & Verified under load (Chapter~\ref{sec:eval}) \\
\bottomrule
\end{tabularx}
\caption{Risks closed.}
\end{table}

\subsubsection{Accepted}

\begin{table}[H]
\centering
\begin{tabularx}{\textwidth}{L{0.36\textwidth} Y}
\toprule
\textbf{Item} & \textbf{Why it is accepted} \\
\midrule
A central-manager restart removes every worker for $\sim$10 minutes & Self-heals; invisible in \kn{condor\_status}. Lower \kn{UPDATE\_INTERVAL} if the window matters \\
Memory eviction is only as fresh as the cron period (20\,s) & A job can allocate hard within that window. The kernel OOM killer is the backstop, and it may pick the wrong victim \\
A job can stay suspended indefinitely on a persistently busy machine & Accepted deliberately: it holds its claim and RSS but no CPU, and resumes when load drops. If it becomes a problem the fix is a cap on suspended time, not a shorter \kn{PREEMPT} \\
Disk-pressure eviction is untested & The expression composes correctly with the entry node's absolute floors, but has not been exercised. Testing it means filling the volume that holds the schedd's queue \\
\bottomrule
\end{tabularx}
\caption{Risks accepted with reasoning.}
\end{table}

\subsubsection{Open}

\begin{table}[H]
\centering
\begin{tabularx}{\textwidth}{L{0.36\textwidth} Y}
\toprule
\textbf{Item} & \textbf{Status} \\
\midrule
Entry node SPOOL shares an 86\%-full volume with the OS, EXECUTE and \file{job\_queue.log} & \textbf{Mitigated, not fixed.} Transfer caps and disk tiers are the whole defence. The structural fix is SPOOL on its own filesystem, where exhausting it degrades the pool instead of destroying the machine \\
The entry node's second interface is now load-bearing & If that interface drops, every worker on that segment loses its return path. Worth monitoring \\
\kn{request\_disk} may be advisory rather than enforced & \textbf{Unverified.} If advisory, one job can fill the disk regardless of what it asked for, and the admission tier is the only protection \\
Entry node is an actively-used desktop owned by another user & \textbf{Risk.} Single point of failure; must be always-on. Give it a DHCP reservation \\
Remote submit with no local Unix account & \textbf{Unverified.} The test user has an account on the entry node, and every submitter so far is that same user \\
Job isolation model (bare versus Apptainer) & \textbf{Undecided.} Blocks growth if the answer is containers \\
Jetson GPU discovery & \textbf{Unverified} \\
Shared storage for Jupyter homes & \textbf{Required, unprovisioned} \\
\kn{MAX\_TRANSFER\_INPUT\_MB} & Set to 2048 as a starting guard; unvalidated against real workloads \\
Tailscale's IPv6 remains in the user PC's \kn{addrs=} list & \textbf{Open, not currently biting.} The interface pin constrains IPv4 only. Nothing has failed, because the entry node advertises and dials IPv4, but a peer preferring IPv6 would dial an unreachable address. \kn{ENABLE\_IPV6 = FALSE} closes it (\S\ref{sec:vpn}) \\
\kn{RANK} as the adoption lever & \textbf{In use}, untested under contention. No second identity has yet competed for that machine \\
Workers lack the name-resolution guard & The entry node has an \kn{ExecStartPre} that waits for \kn{CONDOR\_HOST} to resolve; workers do not. A worker booting before avahi is ready derives the wrong network and idles until the NetworkManager hook restarts Condor \\
Off-LAN machines & mDNS is link-local. Deferred by decision \\
\bottomrule
\end{tabularx}
\caption{Risks still open.}
\end{table}

\subsection{The next three things}

\begin{enumerate}
  \item \textbf{Exercise disk-pressure eviction.} It is the only branch of
        \kn{EVICT\_PRESSURE} never observed firing, and Chapter~\ref{sec:silent}
        is eleven arguments against assuming an unobserved branch works. Fill a
        scratch directory on the flagship worker until the floor is crossed.
        \emph{Deferred by decision, not forgotten.}
  \item \textbf{Give workers the name-resolution guard.} The self-correction path
        exists, but the window is real and lands hardest on the machine most
        likely to move.
  \item \textbf{Decide the isolation model.} It is cheap now and expensive later,
        and it gates how far the pool can grow beyond people who already trust
        each other. Now that the pool runs on colleagues' own desktops rather
        than spare hardware, that is no longer a hypothetical question.
\end{enumerate}
